\scanpage{10}
CHAPITRE VIII

CONSTRUCTION de M\unsure{E|\={E}}SURES INVARIANTES

\handwritten{\S} 1. \uline{Construction d'une mesure de Haar}.

\uline{Dans ce} \handwritten{\S}, $G$ \uline{désignera un groupe localement compact}.

1. \uline{Le théorème d'existence et d'unicité}.

Rappelons (Chap. VII, \handwritten{\S} 2, n\textsuperscript{o}4) qu'on appelle \uline{mesure de Haar} sur $G$ toute mesure positive $\alpha\mathrel{\handwritten{\ne}}0$ sur $G$, invariante à gauche, et qu'une telle mesure est alors une jauge propre à gauche et à droite, satisfaisant à une formule de la forme
\[
d\alpha(axb)=\handwritten{\Delta}_G(b)d\alpha(x)
\]
quels que soient $a\mathrel{\handwritten{\in}}G$, $b\mathrel{\handwritten{\in}}G$, $\handwritten{\Delta}_G$ étant un caractère de $G$ ; on a alors
\[
d\alpha(x^{-1})=\handwritten{\Delta}_G(x^{-1})d\alpha(x)\ .
\]
Si $\handwritten{\Delta}_G=1$, $\alpha$ est invariante à gauche et à droite ; on dit que $\alpha$ est \uline{bi-invariante} ; il en est évidemment ainsi si $G$ est commutatif.

THEOREME 1.- \uline{Sur tout groupe localement compact} $G$, \uline{il existe une mesure de Haar}, \uline{et}, \uline{à un facteur constant près}, \uline{il n'en existe qu'une}.

L'unicité a déjà été démontrée (Chap. VII, \handwritten{\S} 2, n\textsuperscript{o}4, th.1) ; on va démontrer l'existence par une construction explicite. Introduisons d'abord une notation. Soit $\mu=\handwritten{\sum}_{i=1}^{n}a_i\varepsilon_{s_i}$ une mesure de support fini $\handwritten{\{}s_1,\ldots,s_n\handwritten{\}}$ sur $G$ ; alors, pour $f\mathrel{\handwritten{\in}}\handwritten{\unsure{\mathcal{K}|\mathscr{K}}}(G)$, on désignera par $\handwritten{\gamma}(\mu)f$ la fonction $f'\mathrel{\handwritten{\in}}\handwritten{\mathcal{K}}(G)$ définie par
\[
f'(x)=\handwritten{\sum}_{i=1}^{n}a_i f(s_i^{-1}x)\ ;
\]
on notera que, dès qu'on aura construit une mesure de Haar, on aura le droit, en vertu des définitions du Chap. VII, \handwritten{\S} 2, d'écrire $\mu\ast f$ au lieu de $\handwritten{\gamma}(\mu)f$ (cf. la Remarque à la fin du Chap. VII, \handwritten{\S} 2, n\textsuperscript{o}2). L'application $(\mu,f)\mathrel{\handwritten{\longrightarrow}}\handwritten{\gamma}(\mu)f$ ainsi définie est évidemment bilinéaire ; comme on a $\handwritten{\gamma}(\varepsilon_{st})=\handwritten{\gamma}(\varepsilon_s)\handwritten{\gamma}(\varepsilon_t)$, on en conclut qu'on a $\handwritten{\gamma}(\mu\ast\handwritten{\unsure{\nu|v}})=\handwritten{\gamma}(\mu)\handwritten{\gamma}(\handwritten{\nu})$. De plus, $\handwritten{\gamma}(\mu)f$ est $\mathrel{\handwritten{\geq}}0$ si $\mu$ et $f$ sont $\mathrel{\handwritten{\geq}}0$ ; et, si $\mu$ et $f$ sont $\mathrel{\handwritten{\geq}}0$, on ne peut avoir $\handwritten{\gamma}(\mu)f=0$ que si $\mu=0$ et $f=0$.
% archive sheet number: 9

